\documentclass[12pt, a4paper]{ctexart}
\usepackage{fontspec}
\usepackage{xcolor}
\usepackage{hyperref}
\usepackage{geometry}
\usepackage{enumitem}
\usepackage{parskip}
\usepackage{titlesec}
\usepackage{tcolorbox}
\tcbuselibrary{breakable}
\usepackage{setspace}
\usepackage{listings}
\usepackage{amssymb}
\usepackage{tabularx}

% 页面设置
\geometry{left=2.5cm, right=2.5cm, top=2.5cm, bottom=2.5cm}

% 字体设置 (确保系统已安装Times New Roman和SimSun)
\setmainfont{Times New Roman}
\setCJKmainfont{SimSun}

% 超链接设置
\hypersetup{
    colorlinks=true,
    linkcolor=blue!70!black,
    urlcolor=cyan,
    pdftitle={Java程序设计模拟试卷（J卷）- 深度解析讲义},
    pdfauthor={资深Java/JVM教学组}
}

% 颜色定义
\definecolor{myblue}{RGB}{30,100,200}
\definecolor{lightblue}{RGB}{230,240,255}
\definecolor{darkblue}{RGB}{0,50,120}
\definecolor{codebg}{RGB}{245,245,245}

% 解析框样式定义
\newtcolorbox{bluebox}[1][]{
    breakable,
    colback=lightblue,
    colframe=myblue,
    coltitle=darkblue,
    fonttitle=\bfseries,
    title={#1},
    boxrule=0.8pt,
    arc=4pt,
    left=6pt,
    right=6pt,
    top=6pt,
    bottom=6pt,
    before skip=12pt,
    after skip=12pt
}

% 修复跨页环境间距
\usepackage{etoolbox}
\BeforeBeginEnvironment{bluebox}{\par\vfill}%
\AfterEndEnvironment{bluebox}{\par\vfill}%

% 标题格式设置
\titleformat{\section}
  {\normalfont\Large\bfseries\color{myblue}}{\thesection}{1em}{}
\titlespacing*{\section}{0pt}{3.5ex plus 1ex minus .2ex}{1.5ex plus .2ex}

% 代码块样式设置
\lstset{
    language=Java,
    basicstyle=\ttfamily\small,
    backgroundcolor=\color{codebg},
    keywordstyle=\color{blue}\bfseries,
    commentstyle=\color{green!50!black},
    stringstyle=\color{orange},
    numbers=left,
    numberstyle=\tiny\color{gray},
    frame=single,
    breaklines=true,
    columns=fullflexible,
    showstringspaces=false,
    tabsize=4,
    xleftmargin=1.5em,
    xrightmargin=0.5em
}

\begin{document}

\title{\textcolor{myblue}{\textbf{专升本Java程序设计模拟卷-卷11}}}
\author{fifine-专升本}
\date{2026年3月2日}

\maketitle
\thispagestyle{empty}
\newpage

\section{一、单项选择题核心剖析}

\begin{enumerate}[label=\textbf{\arabic*.}]

	\item \textbf{下列哪项是 Java 语言的特性？}
	      \begin{enumerate}[label=\Alph*.]
		      \item 指针运算
		      \item 多重继承
		      \item 自动内存管理
		      \item 预处理器
	      \end{enumerate}

	      \begin{bluebox}{答案与深度解析}
		      \textbf{答案：C. 自动内存管理}

		      \textbf{【核心认知框架】}
		      \begin{itemize}[leftmargin=*, nosep]
			      \item \textbf{题目本质}：考查Java语言设计的底层哲学，即Java与C/C++在内存模型、类结构及编译机制上的核心壁垒。
			      \item \textbf{关键区分点}：引用 vs 指针、单继承+多接口 vs 多继承、GC vs 手动释放。
			      \item \textbf{命题人意图}：筛选出真正理解Java底层隔离机制的考生，而非仅停留在语法背诵层面。
		      \end{itemize}

		      \textbf{【选项原理深度拆解】}

		      \begin{itemize}[leftmargin=*]
			      \item \textbf{A. 指针运算 — 错误（内存安全屏障机制）}
			            \begin{itemize}[leftmargin=*, nosep]
				            \item \textbf{字节码铁证}：
				                  \begin{lstlisting}[language=Java]
Object obj = new Object();
// javap -c 字节码指令：
// 0: new           #2 // class java/lang/Object
// 3: dup
// 4: invokespecial #1 // Method java/lang/Object."<init>":()V
// 7: astore_1         // 存入局部变量表槽位1（引用类型）
\end{lstlisting}
				                  \textbf{JVM保障机制}：Java字节码中只有 \texttt{astore/aload}（引用操作指令），根本不存在类似C语言的 \texttt{p++} 物理内存偏移指令。
				            \item \textbf{HotSpot源码深度解析}：在JVM底层，Java引用被实现为 \texttt{oopDesc}（普通对象指针）。但这个指针被JVM严格接管，开发者无法获取其真实内存绝对地址，彻底杜绝了"野指针"和"内存越界踩踏"的安全隐患。
				            \item \textbf{命题陷阱破解}：
				                  \begin{itemize}
					                  \item 干扰项："Java有NullPointerException，说明有指针"。
					                  \item 真相：异常全称虽带Pointer，但抛出依据是JVMS §6.5中规定的：当 \texttt{aload} 加载的值为 \texttt{null} 且执行实例操作时，JVM自动通过异常表抛出 \texttt{NPE}，并非硬件级别的段错误。
				                  \end{itemize}
			            \end{itemize}

			      \item \textbf{B. 多重继承 — 错误（二义性阻断机制）}
			            \begin{itemize}[leftmargin=*, nosep]
				            \item \textbf{字节码铁证}：
				                  根据JVMS §4.1 \texttt{ClassFile} 结构定义：
				                  \begin{lstlisting}[language=Java]
ClassFile {
    u2             access_flags;
    u2             this_class;
    u2             super_class; // 铁证：这里仅仅是一个 u2 类型，只能指向唯一一个父类！
    u2             interfaces_count;
    u2             interfaces[interfaces_count]; // 接口可以是数组（多重实现）
}
\end{lstlisting}
				            \item \textbf{类加载深度解析}：多重继承会导致"菱形继承（Diamond Problem）"，如果两个父类有相同字段/方法，子类将发生二义性。Java通过限制 \texttt{super\_class} 只能有一个常量池索引，从字节码格式根源上物理锁死了多重继承的可能性。
				            \item \textbf{命题陷阱破解}：
				                  \begin{itemize}
					                  \item 干扰项："Java接口支持多继承，所以Java支持多继承"。
					                  \item 真相：接口的继承调用使用 \texttt{invokeinterface} 指令，类的继承使用 \texttt{invokespecial/invokevirtual}，两者在JVM方法区的方法表搜索算法完全不同。
				                  \end{itemize}
			            \end{itemize}

			      \item \textbf{C. 自动内存管理 — 正确（垃圾收集子系统）}
			            \begin{itemize}[leftmargin=*, nosep]
				            \item \textbf{JVM机制依据}：JVMS §2.5.3 Heap 章节明确规定："The heap is created on virtual machine start-up... Heap storage for objects is reclaimed by an automatic storage management system (known as a garbage collector)."
				            \item \textbf{运行时行为}：通过 \texttt{new} 指令在堆中分配对象后，对象生命周期由GC（如G1/ZGC）的可达性分析算法（GC Roots Tracing）接管。开发者无需（也无法）调用 \texttt{free()} 或 \texttt{delete}。
				            \item \textbf{必考结论}：自动内存管理是Java解决C++内存泄漏（Memory Leak）痛点的最核心特性。
			            \end{itemize}

			      \item \textbf{D. 预处理器 — 错误（编译管线重构）}
			            \begin{itemize}[leftmargin=*, nosep]
				            \item \textbf{编译原理深度解析}：C/C++依赖预处理器在编译前展开 \texttt{\#include} 和 \texttt{\#define}。而Java编译器（Javac）直接在"解析与填充符号表"阶段（Parse and Enter）读取依赖类的元数据，常量则在编译时直接内联入字节码常量池（Constant Value属性），根本不需要独立的预处理阶段。
				            \item \textbf{命题陷阱破解}：
				                  \begin{itemize}
					                  \item 干扰项："import 关键字就是预处理"。
					                  \item 真相：\texttt{import} 仅仅是提供给编译器的"命名空间简写语法糖"，在生成的 \texttt{.class} 文件中，所有类名都会被还原为全限定名（如 \texttt{java/lang/String}），\texttt{import} 语句本身根本不会进入字节码！
				                  \end{itemize}
			            \end{itemize}
		      \end{itemize}

		      \textbf{【应背诵知识点（命题人思维版）】}
		      \begin{enumerate}[leftmargin=*, nosep]
			      \item \textbf{Java与C++核心分水岭速查表}：
			            \begin{tabular}{|c|c|c|c|}
				            \hline
				            \textbf{特性} & \textbf{Java是否支持} & \textbf{底层实现差异}        & \textbf{记忆口诀}                \\
				            \hline
				            指针操作        & \textbf{X}        & 仅暴露安全引用 (\texttt{oop}) & \textcolor{red}{"有引用无指针"}    \\
				            \hline
				            多重继承        & \textbf{X}        & ClassFile结构限单父类        & \textcolor{red}{"单继承多实现"}    \\
				            \hline
				            内存管理        & \checkmark        & 引入GC子系统接管堆区            & \textcolor{green}{"自动回收防泄漏"} \\
				            \hline
				            预处理器        & \textbf{X}        & 语法糖及常量池直接替代            & \textcolor{red}{"无宏无预处理"}    \\
				            \hline
			            \end{tabular}

			      \item \textbf{高频陷阱清单}：
			            \begin{itemize}
				            \item 陷阱1："Java完全消除了内存泄漏" → \textbf{破解}：错！长生命周期对象持有短生命周期对象引用（如ThreadLocal未remove）仍会导致逻辑内存泄漏。
			            \end{itemize}
		      \end{enumerate}

		      \textbf{【命题趋势与实战策略】}
		      \begin{itemize}[leftmargin=*, nosep]
			      \item \textbf{考场秒杀三步法}：
			            \begin{enumerate}
				            \item 看到"特性对比"题型，立刻在脑海中排除具有\textbf{C++浓厚底层色彩}的词汇（指针、宏、预处理、多继承）。
				            \item 锁定具有\textbf{自动化、安全性}色彩的词汇。
				            \item 1秒选出"自动内存管理"或"跨平台（一次编译到处运行）"。
			            \end{enumerate}
		      \end{itemize}
	      \end{bluebox}

\end{enumerate}

\newpage
\section{二、判断题核心剖析}

\begin{enumerate}[label=\textbf{\arabic*.}]
	\item \textbf{Java 虚拟机（JVM）负责将字节码转换为特定平台的机器码。 （ \quad ）}

	      \begin{bluebox}{答案与深度解析}
		      \textbf{答案：\checkmark （正确）}

		      \textbf{【核心认知框架】}
		      \begin{itemize}[leftmargin=*, nosep]
			      \item \textbf{题目本质}：考查JVM的核心职责以及Java实现"Write Once, Run Anywhere (跨平台)"架构的根本原理。
			      \item \textbf{关键区分点}：前端编译（\texttt{.java} $\rightarrow$ \texttt{.class}） vs 后端编译/解释（\texttt{.class} $\rightarrow$ Machine Code）。
			      \item \textbf{命题人意图}：验证考生是否混淆了 \texttt{javac} 编译器和 \texttt{java} 运行时的边界。
		      \end{itemize}

		      \textbf{【命题逻辑深度拆解】}

		      \begin{itemize}[leftmargin=*]
			      \item \textbf{机制深度解析：双层翻译架构 — 正确（JVM规范核心）}
			            \begin{itemize}[leftmargin=*, nosep]
				            \item \textbf{执行引擎铁证}：
				                  \begin{lstlisting}[language=bash]
# 使用 -XX:+PrintCompilation 证明JVM的运行时转换行为
java -XX:+PrintCompilation MyClass
# 输出示例：
#   125   1       3       java.lang.String::hashCode (55 bytes)
#   126   2       3       java.lang.String::equals (81 bytes)
\end{lstlisting}
				                  \textbf{JVM保障机制}：JVM内部包含一个强大的\textbf{执行引擎（Execution Engine）}。其核心任务就是读取方法区中的字节码指令，并将其翻译成当前操作系统（Windows/Linux/Mac）及CPU架构（x86/ARM）能识别的0101机器码。

				            \item \textbf{混合模式（Mixed Mode）深度解析}：
				                  \begin{itemize}
					                  \item \textbf{解释器（Interpreter）}：程序启动时，逐条读取字节码，立刻翻译为机器码执行。优点：启动快；缺点：执行慢。
					                  \item \textbf{即时编译器（JIT - C1/C2/Graal）}：当某段代码执行频率达到阈值（如10000次，变为热点代码），JIT会将其直接编译为高度优化的本地机器码并缓存入CodeCache。
					                  \item JLS/JVMS 明确划分了前端（生成平台无关字节码）和后端（生成平台相关机器码）的职责。题目描述"将字节码转换为机器码"正是JVM后端的精准定义。
				                  \end{itemize}

				            \item \textbf{命题陷阱破解}：
				                  \begin{itemize}
					                  \item 干扰说法："Java代码是由编译器直接编译成机器码的"。
					                  \item \textbf{真相}：Java源代码是被 \texttt{javac} 编译成了**平台无关的字节码**，只有到了运行期，才由 \textbf{JVM} 负责转换为机器码。题目说法完全符合JVM职责边界。
				                  \end{itemize}
			            \end{itemize}
		      \end{itemize}

		      \textbf{【应背诵知识点（命题人思维版）】}
		      \begin{enumerate}[leftmargin=*, nosep]
			      \item \textbf{Java代码生命周期二段论}：
			            \begin{itemize}
				            \item \textbf{前端}：\texttt{Javac} 负责 \texttt{.java} $\rightarrow$ \texttt{.class}（与系统无关，体现"一次编译"）。
				            \item \textbf{后端}：\texttt{JVM执行引擎} 负责 \texttt{.class} $\rightarrow$ 机器码（与系统强相关，体现"到处运行"）。
			            \end{itemize}
			      \item \textbf{高频陷阱清单}：
			            \begin{itemize}
				            \item 陷阱："JVM不仅负责运行代码，也负责将Java源码编译成字节码" → \textbf{破解}：错！编译源码是 \texttt{JDK} 中的 \texttt{javac} 的工作，\texttt{JRE/JVM} 只认字节码不认源码。
			            \end{itemize}
		      \end{enumerate}
	      \end{bluebox}

\end{enumerate}

\newpage
\section{三、填空题核心剖析}

\begin{enumerate}[label=\textbf{\arabic*.}]
	\item[2.] \textbf{在面向对象中，\underline{\hspace{4cm}} 是指一个类具有多种形态，包括方法重载和 \underline{\hspace{4cm}}。}

	      \begin{bluebox}{答案与深度解析}
		      \textbf{答案：多态（Polymorphism） \quad ；\quad 方法重写（Overriding）}

		      \textbf{【核心认知框架】}
		      \begin{itemize}[leftmargin=*, nosep]
			      \item \textbf{题目本质}：考查面向对象三大基石之一"多态"的完整分类体系。
			      \item \textbf{关键区分点}：编译时多态（静态绑定） vs 运行时多态（动态绑定）。
			      \item \textbf{命题人意图}：检查考生是否将多态片面地等同于"父类引用指向子类对象"，而忽略了广义多态的另一半（重载）。
		      \end{itemize}

		      \textbf{【填空项原理深度拆解】}

		      \begin{itemize}[leftmargin=*]
			      \item \textbf{空1：多态 — OOP灵魂机制}
			            \begin{itemize}[leftmargin=*, nosep]
				            \item \textbf{概念深度解析}："多种形态"的术语正是多态。在JVM规范中，多态体现在方法调用的\textbf{分派（Dispatch）}机制上。
			            \end{itemize}

			      \item \textbf{空2：方法重写（Overriding） — 动态绑定的底层逻辑}
			            \begin{itemize}[leftmargin=*, nosep]
				            \item \textbf{字节码铁证}：
				                  \begin{lstlisting}[language=Java]
// 方法重载 (Overloading) 字节码：
// JVM使用 invokestatic / invokespecial，在编译期完全确定调用的具体方法。

// 方法重写 (Overriding) 字节码：
Animal a = new Dog();
a.eat(); 
// javap -c 输出：
// invokevirtual #5 // Method Animal.eat:()V
\end{lstlisting}
				                  \textbf{JVM保障机制（虚方法表 vtable）}：
				                  \begin{itemize}
					                  \item \textbf{重载（Overload）} = \textbf{静态分派（Static Dispatch）}。根据参数的\textbf{静态类型（声明类型）}决定，在 \texttt{javac} 编译期写入字节码。
					                  \item \textbf{重写（Override）} = \textbf{动态分派（Dynamic Dispatch）}。遇到 \texttt{invokevirtual} 指令时，JVM必须在运行期去堆内存中找到对象的\textbf{实际类型（Actual Type）}，查该实际类型的\textbf{虚方法表（vtable）}，将符号引用动态解析为实际的方法入口地址。
				                  \end{itemize}
			            \end{itemize}
		      \end{itemize}

		      \textbf{【应背诵知识点（命题人思维版）】}
		      \begin{enumerate}[leftmargin=*, nosep]
			      \item \textbf{多态双核辨析速查表}：
			            \begin{tabular}{|c|c|c|c|}
				            \hline
				            \textbf{形式} & \textbf{绑定时期} & \textbf{依赖字节码指令}                 & \textbf{决定因素} \\
				            \hline
				            方法重载        & 编译期 (静态)      & \texttt{invokestatic/special}    & 参数静态类型        \\
				            \hline
				            方法重写        & 运行期 (动态)      & \texttt{invokevirtual/interface} & 对象实际类型        \\
				            \hline
			            \end{tabular}
		      \end{enumerate}
	      \end{bluebox}

\end{enumerate}

\newpage
\section{四、程序运行结果题核心剖析}

\begin{enumerate}[label=\textbf{\arabic*.}]
	\item[2.] \textbf{写出下列程序的输出结果：}
	      \begin{lstlisting}[]
class Parent {
    public void method() { System.out.print("Parent "); }
}
class Child extends Parent {
    public void method() { System.out.print("Child "); }
}
public class Test2 {
    public static void main(String[] args) {
        Parent p = new Child();
        p.method();
    }
}
\end{lstlisting}

	      \begin{bluebox}{答案与深度解析}
		      \textbf{答案：Child }

		      \textbf{【核心认知框架】}
		      \begin{itemize}[leftmargin=*, nosep]
			      \item \textbf{题目本质}：经典的多态（动态绑定/方法重写）执行机制考察。
			      \item \textbf{关键区分点}：引用的静态类型（Parent） vs 对象实例的实际类型（Child）在方法调用时的优先级。
			      \item \textbf{命题人意图}：测试对 \texttt{invokevirtual} 动态分派过程的理解。
		      \end{itemize}

		      \textbf{【代码运行原理解析】}
		      \begin{itemize}[leftmargin=*]
			      \item \textbf{字节码执行轨迹追踪}：
			            \begin{lstlisting}[language=Java]
// Test2.main 方法反编译：
0: new           #2 // class Child （在堆区分配Child对象内存）
3: dup
4: invokespecial #3 // 调用Child.<init>()，初始化对象
7: astore_1         // 将引用存入变量 p (静态类型为 Parent)
8: aload_1          // 加载 p 到操作数栈顶
9: invokevirtual #4 // Method Parent.method:()V  <-- 核心指令
\end{lstlisting}

			      \item \textbf{JVM规范（JVMS §6.5）之 invokevirtual 深度解析}：
			            \begin{enumerate}
				            \item \textbf{编译期假象}：在第9行，编译期看到的指令确实是去调用 \texttt{Parent.method()}，因为编译器只能看见变量 \texttt{p} 的声明类型。
				            \item \textbf{运行期真相}：当JVM执行到 \texttt{invokevirtual} 时，它遵循以下寻址算法：
				                  \begin{itemize}
					                  \item \textbf{Step 1}：找到操作数栈顶的引用（即 \texttt{aload\_1} 压入的指针）。
					                  \item \textbf{Step 2}：顺藤摸瓜，沿着指针找到堆（Heap）中真实存在的对象头（Object Header）。
					                  \item \textbf{Step 3}：读取对象头中的类型指针（Klass Pointer），发现这是一个 \texttt{Child} 类的实例。
					                  \item \textbf{Step 4}：查找 \texttt{Child} 类的虚方法表（vtable）。如果 \texttt{Child} 重写了 \texttt{method()}，表中的指针就指向 \texttt{Child.method()}；否则指向父类。
					                  \item \textbf{Step 5}：因为本题重写了该方法，JVM直接将控制权转交 \texttt{Child.method()}。
				                  \end{itemize}
			            \end{enumerate}

			      \item \textbf{命题陷阱破解（变体预警）}：
			            \begin{itemize}
				            \item \textbf{如果调用的是成员变量？}
				                  \begin{lstlisting}[language=Java]
// 若 Parent 有 int x=10; Child 有 int x=20; 
System.out.println(p.x); // 输出 10！
                  \end{lstlisting}
				                  \textbf{真相}：成员变量**没有多态性**！使用的是 \texttt{getfield} 指令，该指令只看静态类型（编译期决定），由于 \texttt{p} 声明为 \texttt{Parent}，因此直接取父类字段。
			            \end{itemize}
		      \end{itemize}

		      \textbf{【命题趋势与实战策略】}
		      \begin{itemize}[leftmargin=*, nosep]
			      \item \textbf{考场秒杀口诀}：
			            \begin{itemize}
				            \item 调方法看\textbf{右边}（对象实际类型， \texttt{new} 的是谁就调谁的重写方法）。
				            \item 调属性看\textbf{左边}（变量声明类型，声明的是谁就取谁的变量）。
			            \end{itemize}
		      \end{itemize}
	      \end{bluebox}

\end{enumerate}

\newpage
\section{五、编程题核心剖析}

\begin{enumerate}[label=\textbf{\arabic*.}]
	\item \textbf{设计一个"几何图形"类及其子类（20分）}\\
	      要求设计抽象类 \texttt{Shape}、子类 \texttt{Circle} 和 \texttt{Rectangle}，并使用对象数组完成多态调用。

	      \begin{bluebox}{答案与深度解析}
		      \textbf{标准答案源码与踩分点拆解}

		      \begin{lstlisting}[language=Java]
// 踩分点1：正确定义抽象类与抽象方法 (4分)
abstract class Shape {
    public abstract double area();
    public abstract double perimeter();
}

// 踩分点2：继承抽象类并封装属性 (6分)
class Circle extends Shape {
    private double radius; // 必须私有
    public Circle(double radius) { this.radius = radius; }
    
    @Override // 必须实现抽象方法
    public double area() { return Math.PI * radius * radius; }
    
    @Override
    public double perimeter() { return 2 * Math.PI * radius; }
}

// 踩分点3：另一个完整子类的实现 (6分)
class Rectangle extends Shape {
    private double length, width;
    public Rectangle(double length, double width) {
        this.length = length; this.width = width;
    }
    
    @Override
    public double area() { return length * width; }
    
    @Override
    public double perimeter() { return 2 * (length + width); }
}

// 踩分点4：对象数组与多态遍历测试 (4分)
public class Main {
    public static void main(String[] args) {
        // 多态：父类数组存放子类对象
        Shape[] shapes = new Shape[4];
        shapes[0] = new Circle(5.0);
        shapes[1] = new Circle(3.0);
        shapes[2] = new Rectangle(4.0, 5.0);
        shapes[3] = new Rectangle(2.0, 3.0);
        
        // 动态绑定执行对应方法
        for (Shape shape : shapes) {
            System.out.println("面积: " + shape.area() + 
                               ", 周长: " + shape.perimeter());
        }
    }
}
\end{lstlisting}

		      \textbf{【JVM层面的架构设计深度解析】}
		      \begin{itemize}[leftmargin=*, nosep]
			      \item \textbf{抽象类的字节码本质（ACC\_ABSTRACT）}：
			            当声明 \texttt{abstract class Shape} 时，JVM会在 \texttt{ClassFile} 结构的 \texttt{access\_flags} 中打上 \texttt{ACC\_ABSTRACT} 标记。
			            一旦类有此标记，JVM规范第5.3节强制规定：执行 \texttt{new Shape()} 指令时，类加载器会在校验阶段直接抛出 \texttt{InstantiationError}，从物理层面保证了模板类的不可实例化。

			      \item \textbf{对象数组的内存布局（ArrayKlass）}：
			            \texttt{new Shape[4]} 生成的并不是4个图形对象，而是一个类型为 \texttt{[LShape;} 的数组引用，堆中分配的是4个 \texttt{null} 槽位（占用 4 * 8 = 32字节，64位指针）。
			            直到 \texttt{shapes[0] = new Circle()}，堆内存才真正创建实例，并将引用存入数组。

			      \item \textbf{Foreach遍历的多态分派}：
			            当调用 \texttt{shape.area()} 时，虽然 \texttt{shape} 引用静态类型是 \texttt{Shape}，但由于每个数组元素存的实际是 \texttt{Circle} 或 \texttt{Rectangle} 的句柄，JVM依然通过 \texttt{invokevirtual} 指令，完美实现"同一段代码，不同对象表现出不同行为"的终极多态。
		      \end{itemize}

		      \textbf{【判卷人视角的扣分陷阱】}
		      \begin{itemize}[leftmargin=*, nosep]
			      \item \textbf{致命陷阱1}：子类属性没有使用 \texttt{private} 修饰。破坏了三大特性之一的"封装"。
			      \item \textbf{致命陷阱2}：忘记写构造方法。由于存在带参需求，如果未显式提供带参构造器，编译器默认提供无参构造器，导致对象属性无法初始化。
			      \item \textbf{致命陷阱3}：将测试类写死了四个对象变量，而不是使用 \texttt{Shape[]} 数组循环。此举直接导致多态核心能力得不到展示，最高扣4分。
		      \end{itemize}
	      \end{bluebox}

\end{enumerate}

\end{document}